\subsection{Qi et al.\ (2023): MTP for L1$_0$-TiAl and D0$_{19}$-Ti$_3$Al}
\label{sec:appendix-example-qi2023}

\paragraph{Q1 -- Bibliographic identification.}
Qi, Aitken, Pei, Tan, Zuo, Jhon, Quek, Wen, Wu, and Ong train a single
Moment Tensor Potential covering the binary Ti--Al system, with focus on
the dual-phase $\gamma$-TiAl (L1$_0$) and $\alpha_2$-Ti$_3$Al
(D0$_{19}$) intermetallics. Published in \emph{Phys.\ Rev.\ Materials}
7, 103602 (2023); arXiv:2305.11825.
\lstinputlisting[language=turtle]{artifacts/kg/listings/qi2023/q01-bibliographic.ttl}

\paragraph{Q2 -- Material system.}
The MTP is trained over the Ti--Al binary phase diagram and benchmarked
on the two technologically relevant intermetallics, $\gamma$-TiAl
(L1$_0$, P4/mmm) and $\alpha_2$-Ti$_3$Al (D0$_{19}$, P6$_3$/mmc), as
well as on FCC-Al, HCP-Ti, TiAl$_2$, and TiAl$_3$. We model this as
one $\MaterialSystemC$ at the binary level, with the phase information
recorded in $\dlRole{materialClass}$.
\lstinputlisting[language=turtle]{artifacts/kg/listings/qi2023/q02-material.ttl}

\paragraph{Q3 -- Reference calculation method.}
Reference data are produced from spin-polarised DFT calculations.
\lstinputlisting[language=turtle]{artifacts/kg/listings/qi2023/q03-reference-method-type.ttl}

\paragraph{Q4 -- Reference settings.}
DFT calculations use VASP with the PBE/GGA exchange--correlation
functional and PAW pseudopotentials (3p$^6$3d$^3$4s$^1$ valence for
Ti, 3s$^2$3p$^1$ for Al), a 520~eV plane-wave cutoff, and
Monkhorst--Pack k-point meshes with density $\geq 100/\text{\AA}^3$
(a single $\Gamma$-point is used during AIMD). Frozen-core treatment
beyond the PAW partition is not explicitly reported.
\lstinputlisting[language=turtle]{artifacts/kg/listings/qi2023/q04-reference-settings.ttl}

\paragraph{Q5 -- Training dataset.}
The training set comprises 3{,}798 configurations produced from five
sampling categories (see Q6): 33 ground-state polymorphs, 2{,}160
AIMD snapshots, 185 surface structures, 580 solid-solution
configurations, and 840 strained supercells. Configurations cover
energies and forces; stresses appear in the loss function with weight
zero, so they are not effectively learned and not recorded as a
covered property. A 90:10 train:test split is applied.
\lstinputlisting[language=turtle]{artifacts/kg/listings/qi2023/q05-dataset.ttl}

\paragraph{Q6 -- Sampling strategies.}
Five strategies were combined: ground-state polymorph enumeration
from the Materials Project, vibrational sampling via NVT AIMD at 300,
1000, and 3000~K and at 90/100/110\% of the ground-state volume,
surface enumeration with Miller indices up to three, solid-solution
chemical sampling at 12.5\% steps, and homogeneous strain sampling
($\pm10\%$ in 1\% intervals across six modes).
\lstinputlisting[language=turtle]{artifacts/kg/listings/qi2023/q06-sampling.ttl}

\paragraph{Q7 -- MLIP method, components, implementation.}
The method is a Moment Tensor Potential (MTP) implemented in the MLIP
package~\cite{shapeev2016mtp,novikov2021mlip} for training,
with downstream simulations in LAMMPS via the maml workflow library.
The loss function is a weighted MSE with energy:force:stress weights
of $100{:}1{:}0$. The training algorithm is not explicitly named in
the paper.
\lstinputlisting[language=turtle]{artifacts/kg/listings/qi2023/q07-method.ttl}

\paragraph{Q8 -- Hyperparameter settings.}
A grid search over cutoff radius $r_c \in [4.4, 7.0]$\,\AA{} (steps of
0.2\,\AA) and maximum level $\mathit{lev}_{\max} \in [18, 24]$ (even
integers), with five random initialisations per cell, was performed
(280 candidate MTPs in total). The selected production MTP uses
$r_c = 4.8$\,\AA{} and $\mathit{lev}_{\max} = 22$. Other implementation
hyperparameters (radial-basis count, regularisation, optimiser
schedule) are not reported.
\lstinputlisting[language=turtle]{artifacts/kg/listings/qi2023/q08-hyperparameter-settings.ttl}

\paragraph{Q9 -- MLIP run and trained model.}
A single training run produces one trained MTP that is used across
all benchmarked phases. No active-learning or fold-ensemble structure
is reported.
\lstinputlisting[language=turtle]{artifacts/kg/listings/qi2023/q09-run-and-model.ttl}

\paragraph{Q10 -- Benchmark results.}
Headline accuracy figures are MAEs of formation energies of (i)
surface structures (8\,meV/atom for the MTP, vs.\ 51\,meV/atom for
the MLP3 of Seko) and (ii) ground-state polymorphs (18\,meV/atom for
the MTP, vs.\ 13\,meV/atom for MLP3). The MTP further attains
$\Delta E_{\mathrm{EOS}} < 2$\,meV/atom in the L1$_0$-TiAl,
D0$_{19}$-Ti$_3$Al, HCP-Ti, and FCC-Al phases. Per-property numerical
errors on elastic constants, surface energies, and GSFEs are
discussed at length in the paper but reported as derived figures
rather than as a single accuracy metric per property.
\lstinputlisting[language=turtle]{artifacts/kg/listings/qi2023/q10-benchmark-results.ttl}

\paragraph{Q11 -- Computational resources.}
\emph{Not reported.} The paper gives no training duration, GPU hours,
training hardware, peak memory, inference time per atom, or
inference hardware.
\lstinputlisting[language=turtle]{artifacts/kg/listings/qi2023/q11-resources.ttl}

\paragraph{Q12 -- Gaps.}
Beyond Q11 (compute resources entirely absent), the following
ontology-expressible items are not reported: the training algorithm
or optimiser used by the MLIP package; the MLIP package version;
detailed MTP architectural hyperparameters beyond $r_c$ and
$\mathit{lev}_{\max}$ (radial/angular basis counts, regularisation,
loss-function temperature scheduling); explicit
$\dlRole{frozenCore}$ treatment beyond the PAW partition; and a
controlled-vocabulary individual for AIMD as a sampling strategy
(currently encoded as an ad-hoc instance with \texttt{rdfs:label} and
\texttt{rdfs:comment}).
